\documentclass[11pt,a4paper]{article}
\usepackage[utf8]{inputenc}
\usepackage[vietnamese]{babel}
\usepackage{amsmath,amssymb}
\usepackage{geometry}
\geometry{a4paper, total={180mm,260mm}, left=15mm, right=15mm, top=18mm, bottom=18mm}
\usepackage{tikz}
\usetikzlibrary{shapes.geometric, arrows.meta, positioning, calc}
\usepackage{tcolorbox}
\tcbuselibrary{skins,breakable}
\usepackage{xcolor}
\usepackage{multicol}
\usepackage{array}
\usepackage{fontawesome}
\usepackage{fancyhdr}

\pagestyle{fancy}
\fancyhf{}
\renewcommand{\headrulewidth}{0.5pt}
\renewcommand{\footrulewidth}{0.5pt}
\fancyhead[L]{\small\bfseries\color{teal}\faBook\ SỔ TAY KIẾN THỨC TOÁN 3}
\fancyhead[R]{\small\bfseries\color{teal}Lớp Toán Thầy Lĩnh -- loptoanthaylinh.vn}
\fancyfoot[C]{\small\color{gray}-- Sổ tay kiến thức: Bài 40. Luyện tập chung: Phép nhân, phép chia trong phạm vi 1 000 --}

\definecolor{primary}{HTML}{005F73}
\definecolor{secondary}{HTML}{0A9396}
\definecolor{accent}{HTML}{EE9B00}
\definecolor{danger}{HTML}{AE2012}
\definecolor{boxbg}{HTML}{F4F6F8}

\begin{document}

% --- BANNER TIÊU ĐỀ ---
\begin{tcolorbox}[
  enhanced,
  colback=white,
  colframe=primary,
  arc=3mm,
  boxrule=1.2pt,
  left=4mm, right=4mm, top=2.5mm, bottom=2.5mm,
  drop shadow=black!10!white
]
  \centering
  {\Large\bfseries\color{primary}\faBookmark\ BÀI 40: LUYỆN TẬP CHUNG: NHÂN CHIA PHẠM VI 1 000}\\[1.5mm]
  {\small\itshape\color{black!75}Tổng kết kĩ thuật tính nhân chia 3 chữ số, so sánh số lần và bài toán hai bước trên 1 trang A4}
\end{tcolorbox}

\vspace{-2mm}


\noindent
\begin{minipage}[t]{0.49\textwidth}\vspace{0pt}
\begin{tcolorbox}[
  enhanced,
  colback=blue!4!white,
  colframe=blue!75!black,
  arc=2mm, boxrule=0.8pt,
  left=2mm, right=2mm, top=1.8mm, bottom=1.8mm,
  title={\footnotesize\bfseries\color{white}\faCalculator\ 1. Kĩ thuật nhân chia phạm vi 1 000}
]
{\footnotesize
\begin{itemize}\setlength{\itemsep}{1.5pt}\setlength{\leftskip}{-2.5mm}
  \item \textbf{Nhân $3$ chữ số với $1$ chữ số:}
  Nhân lần lượt từ phải sang trái (đơn vị $\to$ chục $\to$ trăm), nhớ sang hàng liền trước nếu tích $\ge 10$.
  \item \textbf{Chia $3$ chữ số cho $1$ chữ số:}
  Chia lần lượt từ trái sang phải theo vòng 3 nhịp: Chia $\to$ Nhân $\to$ Trừ.
  \item Khi lượt chia bé hơn số chia $\implies$ ghi chữ số $\mathbf{0}$ vào thương rồi mới hạ tiếp.
\end{itemize}
}
\end{tcolorbox}
\end{minipage}
\hfill
\begin{minipage}[t]{0.49\textwidth}\vspace{0pt}
\begin{tcolorbox}[
  enhanced,
  colback=teal!4!white,
  colframe=teal!75!black,
  arc=2mm, boxrule=0.8pt,
  left=2mm, right=2mm, top=1.8mm, bottom=1.8mm,
  title={\footnotesize\bfseries\color{white}\faBalanceScale\ 2. Kĩ năng so sánh số lần}
]
{\footnotesize
\begin{itemize}\setlength{\itemsep}{1.5pt}\setlength{\leftskip}{-2.5mm}
  \item \textbf{Số lớn gấp mấy lần số bé:}
  $$\text{Số lần} = \text{Số lớn} : \text{Số bé}$$
  \item \textbf{Số bé bằng một phần mấy số lớn:}
  $$\text{Nếu } a : b = k \implies b = \mathbf{\dfrac{1}{k}} \text{ của } a$$
  \item \textbf{Giải toán bằng hai bước tính:}
  Tìm ẩn số trung gian ở bước 1, trả lời câu hỏi tổng thể ở bước 2.
\end{itemize}
}
\end{tcolorbox}
\end{minipage}

\vspace{-2.5mm}
% --- KHỐI 3: SƠ ĐỒ TƯ DUY TIKZ ---
\noindent
\begin{tcolorbox}[
  enhanced,
  colback=orange!4!white,
  colframe=orange!80!black,
  arc=2mm, boxrule=0.8pt,
  left=2mm, right=2mm, top=1.5mm, bottom=1.5mm,
  title={\footnotesize\bfseries\color{white}\faSitemap\ SƠ ĐỒ TƯ DUY: HỆ THỐNG PHÉP TÍNH VÀ SO SÁNH LỚP 3}
]
\centering
\begin{tikzpicture}[
  box/.style={draw, rounded corners=2pt, align=center, font=\scriptsize\bfseries, inner sep=2.5pt, minimum height=0.75cm},
  arr/.style={-{Stealth[length=2mm]}, thick}
]
  \node[box, fill=purple!15, draw=purple!80!black, text width=2.6cm] (p1) at (0, 0) {ĐẶT TÍNH NHÂN\\Từ phải sang trái\\Nhân trước cộng nhớ};
  \node[box, fill=teal!15, draw=teal!80!black, text width=2.6cm] (p2) at (3.5, 0) {ĐẶT TÍNH CHIA\\Từ trái sang phải\\Chia - Nhân - Trừ};
  \node[box, fill=blue!12, draw=blue!80!black, text width=2.6cm] (p3) at (7.0, 0) {SO SÁNH GẤP LẦN\\Lấy Lớn $:$ Bé\\Đơn vị: lần};
  \node[box, fill=orange!15, draw=orange!80!black, text width=2.6cm] (p4) at (10.5, 0) {GIẢI BÀI TOÁN\\2 bước tính logic\\Đọc kĩ câu hỏi cuối};

  \draw[arr, purple!80!black] (p1) -- (p2);
  \draw[arr, teal!80!black] (p2) -- (p3);
  \draw[arr, blue!80!black] (p3) -- (p4);
\end{tikzpicture}
\end{tcolorbox}

\vspace{-2.5mm}
% --- KHỐI 4: VÍ DỤ & MẸO TRÁNH BẪY ---
\noindent
\begin{minipage}[t]{0.49\textwidth}\vspace{0pt}
\begin{tcolorbox}[
  enhanced,
  colback=green!4!white,
  colframe=green!60!black,
  arc=2mm, boxrule=0.8pt,
  left=2mm, right=2mm, top=1.8mm, bottom=1.8mm,
  title={\footnotesize\bfseries\color{white}\faCheckCircle\ Ví dụ thực hành chuẩn mực}
]
{\footnotesize
\textbf{1. Đặt tính rồi tính:}\\
$154 \times 3 = \mathbf{462}$ (nhớ 1 sang chục).\\
$618 : 3 = \mathbf{206}$ (chữ số 0 ở hàng chục của thương).\\
\textbf{2. Bài toán so sánh:}\\
Bao gạo to nặng $45\text{ kg}$, bao gạo nhỏ nặng $9\text{ kg}$.\\
-- Bao to gấp bao nhỏ: $45 : 9 = \mathbf{5}$ (lần).\\
-- Bao nhỏ bằng $\mathbf{\dfrac{1}{5}}$ bao to.
}
\end{tcolorbox}
\end{minipage}
\hfill
\begin{minipage}[t]{0.49\textwidth}\vspace{0pt}
\begin{tcolorbox}[
  enhanced,
  colback=red!4!white,
  colframe=red!70!black,
  arc=2mm, boxrule=0.8pt,
  left=2mm, right=2mm, top=1.8mm, bottom=1.8mm,
  title={\footnotesize\bfseries\color{white}\faExclamationTriangle\ Mẹo tránh bẫy trắc nghiệm}
]
{\footnotesize
\begin{itemize}\setlength{\itemsep}{1.5pt}\setlength{\leftskip}{-2.5mm}
  \item \textbf{Bẫy quên chữ số 0 ở thương khi chia:}\\
  $618 : 3$, lấy $6 : 3 = 2$; hạ $1 < 3$ phải ghi ngay \textbf{0 vào thương} rồi mới hạ $8$. Viết thiếu số 0 thành $26$ là sai nghiêm trọng!
  \item \textbf{Bẫy nhầm giữa trừ và chia:}\\
  Hỏi "gấp mấy lần" $\implies$ Chia.\\
  Hỏi "hơn kém mấy kg" $\implies$ Trừ.
\end{itemize}
}
\end{tcolorbox}
\end{minipage}


\end{document}
